% =====================================================================
%  CV Harshil Gorasiya (English)
%  Kompilieren mit pdflatex (zweimal). Legen Sie photo.png in denselben Ordner.
%  Datei fuer den Versand umbenennen: Lebenslauf_Harshil_Gorasiya.pdf
%  Layout bewusst ATS-freundlich: Standardschrift, Klartext, keine Icons,
%  keine Fortschrittsbalken, einspaltiger Textfluss, eindeutige Ueberschriften.
% =====================================================================
\documentclass[10pt,a4paper]{article}

\usepackage[a4paper,left=1.7cm,right=1.7cm,top=1.1cm,bottom=1.35cm]{geometry}
\usepackage[T1]{fontenc}
\usepackage[utf8]{inputenc}
\usepackage{textcomp}
\usepackage[scaled=0.96]{helvet}
\renewcommand{\familydefault}{\sfdefault}
\usepackage[english]{babel}
\usepackage{microtype}
\usepackage{xcolor}
\usepackage{graphicx}
\usepackage{enumitem}
\usepackage{titlesec}
\usepackage{array}
\usepackage{fancyhdr}
\usepackage{lastpage}
\makeatletter

\definecolor{navy}{HTML}{1B365D}
\definecolor{mid}{HTML}{5A6472}
\definecolor{tint}{HTML}{EEF2F7}

\usepackage[unicode,colorlinks=true,urlcolor=navy,linkcolor=navy,
  pdftitle={CV Harshil Gorasiya},
  pdfauthor={Harshil Gorasiya},
  pdfsubject={CV: AI/ML Engineering, LLM, RAG, AI agents},
  pdfkeywords={working student, master thesis, AI, machine learning, LLM, RAG, Python, PyTorch}]{hyperref}

\input{glyphtounicode}
\pdfgentounicode=1

% Layout bleibt in jeder Umgebung identisch (keine Silbentrennung, Flattersatz)
\hyphenpenalty=10000
\exhyphenpenalty=10000
\emergencystretch=3em
\setlength{\parindent}{0pt}
\setlength{\parskip}{0pt}
\raggedright
\newlength{\leftcol}\setlength{\leftcol}{2.45cm}
\newcommand{\gap}{0.5em}
\setlist[itemize]{leftmargin=*,labelindent=\dimexpr\leftcol+\gap\relax,labelwidth=0.55em,labelsep=0.4em,align=left,itemsep=0.6pt,topsep=1.5pt,parsep=0pt,partopsep=0pt,label={\color{navy}\textbullet}}


% Fusszeile
\pagestyle{fancy}
\fancyhf{}
\renewcommand{\headrulewidth}{0pt}
\renewcommand{\footrulewidth}{0pt}
\fancyfoot[C]{\footnotesize\color{mid}Harshil Gorasiya \textperiodcentered{} CV \textperiodcentered{} Page \thepage{} of \pageref{LastPage}}

% Abschnittsueberschriften
\titleformat{\section}{\normalfont\bfseries\color{navy}\normalsize}{}{0pt}
  {\MakeUppercase}[\vspace{-0.45em}{\color{navy}\rule{\linewidth}{0.7pt}}]
\titlespacing*{\section}{0pt}{9.5pt}{5pt}

% Hilfsmakros
\newcommand{\sep}{\hspace{0.45em}\textbar\hspace{0.45em}}
\newcommand{\dotsep}{\hspace{0.3em}\textperiodcentered\hspace{0.3em}}
\newcommand{\entry}[3]{%
  \begingroup\raggedright\leftskip=\dimexpr\leftcol+\gap\relax\noindent
  \hspace*{-\dimexpr\leftcol+\gap\relax}%
  \vtop to 0pt{\hsize=\dimexpr\leftcol+\gap\relax\raggedright\small\color{mid}\noindent #1\par\vss}%
  #2\par #3\par\endgroup\vspace{4pt}}
\newcommand{\etitle}[2]{\textbf{#1}\ifx&#2&\else\textcolor{mid}{\sep}#2\fi}
\newcommand{\stack}[1]{{\small\itshape\color{mid}#1}\par\vspace{1pt}}
\newcommand{\kpi}[2]{\begin{minipage}[t]{\linewidth}\centering{\large\bfseries\color{navy}#1}\\[-1pt]{\scriptsize\color{mid}#2}\end{minipage}}

\renewcommand\normalsize{\@setfontsize\normalsize{9.6}{11.3}}
\makeatother
\begin{document}
\normalsize

\noindent
\begin{minipage}[t]{\dimexpr\textwidth-3.05cm\relax}
  {\fontsize{25}{28}\selectfont\bfseries\color{navy}Harshil Gorasiya}\\[4pt]
  {\bfseries LLM systems: retrieval, agents, and small-model research}\\[1pt]
  {\color{navy}\bfseries LLM Applications\dotsep RAG\dotsep AI Agents\dotsep LLM Research}\\[5pt]
  {\small
  Heilbronn, Germany\sep +49 155 6351 7346\sep
  \href{mailto:harshil.gorasiya.0011@gmail.com}{harshil.gorasiya.0011@gmail.com}\\[1pt]
  LinkedIn: \href{https://www.linkedin.com/in/harshil-gorasiya-054a58319}{linkedin.com/in/harshil-gorasiya-054a58319}\\[1pt]
  GitHub: \href{https://github.com/N0t-Harshil}{github.com/N0t-Harshil}\sep
  Portfolio: \href{https://harshil-0011.github.io/My-Portfolio/}{harshil-0011.github.io/My-Portfolio}}
\end{minipage}\hfill
\begin{minipage}[t]{2.7cm}
  \raggedleft
  \IfFileExists{photo.png}{\setlength{\fboxsep}{0pt}\setlength{\fboxrule}{0.6pt}%
    \fcolorbox{navy}{white}{\includegraphics[width=2.7cm]{photo.png}}}{}
\end{minipage}

\vspace{7pt}
\noindent\colorbox{tint}{\parbox{\dimexpr\textwidth-2\fboxsep\relax}{\small
  \textbf{Seeking:} Working student position (20 hrs/week during semester)\dotsep
  Full-time during semester breaks\dotsep LLM/RAG role after graduation (2027)\\[1pt]
  \textbf{Available:} immediately\sep
  \textbf{Residence status:} valid residence permit with work permit}}

\vspace{6pt}
\noindent
\begin{tabular*}{\textwidth}{@{\extracolsep{\fill}}*{5}{>{\centering\arraybackslash}p{0.185\textwidth}}}
  \kpi{0.654}{nDCG@10, BEIR SciFact, 300 queries} &
  \kpi{8.1 ms}{overhead per agent step, laptop CPU} &
  \kpi{28--32 FPS}{YOLOv8n on CPU only} &
  \kpi{79\,\%}{accuracy (up from 71\,\%)} &
  \kpi{250 tests}{pytest suite, Ekdant}
\end{tabular*}

\section*{Profile}
Master's student at Heilbronn University focused on LLM applications. I build local RAG systems with hybrid retrieval and AI agents with tool use, and I research a reasoning module for language models in a supervised project. I also bring two ML internships and a role in AI-enabled web development.

\section*{Education}
\entry{since 03/2025}
  {\etitle{M.Sc. Software Engineering and Management}{Hochschule Heilbronn}}
  {\small English-taught programme\sep GPA so far: 1.9 (German scale)\sep expected graduation: 2027}
\entry{2020 -- 2024}
  {\etitle{B.Tech. Information Technology}{A.\,D. Patel Institute of Technology, Gujarat}}
  {\small India\sep Grade: 1.9 (converted to German scale)}

\section*{Projects}
\entry{AI agent}
  {\etitle{Ardan-CLI: Autonomous Terminal Coding Agent}{}}
  {\stack{TypeScript (strict), Node.js 24, 0 runtime dependencies, MCP, LSP}
  \begin{itemize}
    \item Agent loop with streaming, \textbf{35 tools}, plan mode, sub-agents, and checkpoints; \textbf{8 LLM providers} with automatic failover.
    \item SSE/NDJSON streaming, MCP and LSP clients without SDKs; fallback parser for models without function calling.
    \item \textbf{8.1 ms} overhead per agent step (laptop CPU); CI on Linux, macOS, Windows.
  \end{itemize}}
\entry{Research engine}
  {\etitle{Local Perplex: Local AI Research Assistant with Citations}{}}
  {\stack{Python, FastAPI, SSE, SQLite, Ollama, Docker}
  \begin{itemize}
    \item Searches the web and your own documents (PDF, Word, Excel, and more), reads pages in parallel, and streams cited answers; \textbf{26 LLM providers}, 7 search backends, fully local if desired.
    \item Own BM25 ranker: \textbf{0.654 nDCG@10} on BEIR SciFact, 300 queries (previous version 0.038; Anserini reference 0.665), 9.5 ms per query.
    \item Answer page in about \textbf{200 ms}; SQLite knowledge graph (77\texttimes{} faster writes); \textbf{55 tests}, 76\,\% coverage, CI on Ubuntu and Windows.
  \end{itemize}}
\entry{RAG system}
  {\etitle{ANVITHA: Local Knowledge Vault with Grounded, Cited Answers}{}}
  {\stack{Python, FastAPI, Ollama, SQLite, NumPy, BM25}
  \begin{itemize}
    \item Turns PDF, TXT, and Markdown files into a linked Markdown knowledge vault; answers come only from the vault and cite the original page.
    \item Facts are checked against the source text before they are stored; every answer sentence is checked against the vault afterwards, and unsupported parts are removed.
    \item Questions outside the vault are refused in under 0.3 s without a model call.
    \item Tested with a local 0.8B-parameter model (Ollama): 8 of 8 covered questions answered (12 test questions in total), median 7.4 s to full answer; 26 automated tests in about 1.5 s.
  \end{itemize}}
\newpage
\entry{Master's thesis\\(ongoing)}
  {\etitle{Ekdant: Debate-Sandbox Reasoning Module for Language Models}{}}
  {\stack{Hochschule Heilbronn, supervised by a professor\sep Python, PyTorch, Hugging Face}
  \begin{itemize}
    \item Research question: does a parallel debate module (4 trajectories, 2 debate rounds, consensus) improve a language model's reasoning? Implemented in pure PyTorch as a zero-initialized residual module (\textbf{+19.59M parameters}) on a pretrained SmolLM2-135M.
    \item Controlled A/B comparison against an architecture-matched baseline on generative medical QA (PubMedQA, MedQuAD); leakage-free split by question; evaluated with perplexity, ROUGE-L, token F1, and paired-bootstrap confidence intervals.
    \item Best variant (module as final stage) beats the baseline on every metric, but only narrowly (perplexity \textbf{3.783 to 3.776}, one run); the from-scratch pilot showed no difference. Second seed planned to confirm.
    \item \textbf{250 pytest tests}; found and fixed 4 design defects; literature base of 66 papers.
  \end{itemize}}
\entry{More projects}
  {\etitle{AI Video Upscaling}{Python, PyTorch}}
  {\small Windows desktop application for video super-resolution with model selection and runtime monitoring; compared Real-ESRGAN variants using PSNR and SSIM.}
\entry{\strut}
  {\etitle{MNIST-MLOps}{Python, PyTorch, Flask, SQLite}}
  {\small Demo of an MLOps workflow: incremental training, metrics and model versions in SQLite, web dashboard.}

\section*{Professional Experience}
\entry{05/2024 -- 02/2025\\(10 months)}
  {\etitle{AI \& Web Solutions Developer}{Aani Brothers Infotech, Surat, India}}
  {\begin{itemize}
    \item Built web applications that translate business requirements into user interfaces and application workflows.
    \item Contributed to web projects with AI-enabled features; integrated functional requirements into the user experience.
    \item Conducted user research, wireframing, and prototyping; iterated on designs to improve usability and consistency.
  \end{itemize}}
\entry{12/2023 -- 04/2024\\(5 months)}
  {\etitle{Deep Learning Intern}{PD Avenue, Surat, India}}
  {\begin{itemize}
    \item Built a residual-block CNN (TensorFlow/Keras) for age and gender classification on \textbf{22,000+ images}; raised validation accuracy from \textbf{71\,\% to 79\,\%} (architecture optimization, class-balanced augmentation).
    \item Cut training time on V100 GPUs by \textbf{20\,\%} (cosine learning-rate scheduling, larger batch size) with no loss of accuracy.
    \item Found and fixed a data-leakage issue; ensured reliability across \textbf{10+ structured training runs}.
  \end{itemize}}
\entry{05/2023 -- 06/2023\\(2 months)}
  {\etitle{Machine Learning Intern}{Dotcom IoT LLP, Surat, India}}
  {\begin{itemize}
    \item YOLOv8n object detection on a CPU-only edge device (Intel Core i7): 720p video at \textbf{28--32 FPS}; \textbf{2.4\texttimes{} higher throughput} than the YOLOv5s baseline.
    \item Curated and annotated \textbf{3,000+ images} across five vehicle classes.
    \item Tuned IoU and NMS parameters: reduced false positives in low light by \textbf{18\,\%} (dedicated night-test dataset).
    \item Packaged inference and visualization into a deployment-ready Python application with performance benchmarking.
  \end{itemize}}

\section*{Skills}
\entry{LLM \& RAG}
  {\small RAG pipelines, hybrid retrieval (BM25, embeddings), grounded answering, Ollama, ReAct agents and tool use, Anthropic API, Model Context Protocol (MCP), prompt engineering}{}
\entry{LLM research}
  {\small Fine-tuning, Transformer architectures (GQA, SwiGLU), A/B experiment design, PyTorch, Hugging Face Transformers}{}
\entry{ML \& Vision}
  {\small PyTorch, TensorFlow/Keras, scikit-learn, YOLOv8, CNNs, OpenCV}{}
\entry{Engineering}
  {\small Python, TypeScript, SQL, Bash, Git, Docker, Linux, pytest, CI (GitHub Actions), REST APIs, FastAPI, Flask, MLflow, SQLite, PostgreSQL, MongoDB, AWS (basic)}{}
\entry{Evaluation}
  {\small Latency and throughput benchmarking, nDCG@10 (BEIR), validation accuracy, false-positive analysis, PSNR/SSIM, perplexity, ROUGE-L, token F1, paired bootstrap}{}

\section*{Certifications}
\entry{2026}{\small Claude with the Anthropic API\sep Model Context Protocol: Advanced Topics (both Anthropic)}{}
\entry{2025}{\small Machine Learning with Python (IBM)\sep Fundamentals of Data Science (Databricks)}{}
\entry{Further}{\small CRISP-DM on the AWS Stack\sep Big Data with Spark and Hadoop\sep AWS Academy Cloud}{}

\section*{Languages}
\entry{\strut}{\small \textbf{English:} fluent\sep \textbf{German:} A2--B1 (self-assessed)\sep \textbf{Hindi:} fluent\sep \textbf{Gujarati:} native}{}

\vspace{2pt}
{\small\color{mid}Heilbronn, 05.11.2026}

\end{document}
